\chapter{Simulation Design and Data}
\label{ch:simulation}

\section{Overall Research Design}
\label{sec:design}

The thesis uses a two-part mixed-method design. Part~I,
\Cref{ch:design}, is a conceptual and architectural analysis answering
RQ1 by theoretical synthesis supported by literature. Part~II, this
chapter and the next, is a quantitative simulation study answering RQ2.

The two parts are linked in both directions. The design blocks derived
in Part~I define the parameters that Part~II varies: the batch interval
sweep operationalises \Cref{sec:interval}, and the rationing rule
comparison operationalises \Cref{sec:rationing}. In the other direction,
the measured results feed into the architectural recommendation of
\Cref{ch:discussion}, and the predictions of \Cref{tab:predictions} that
the simulation cannot test define the boundary of what that
recommendation can claim.

The identifying idea of Part~II is simple. A single historical order
flow is replayed twice, once through a continuous engine and once
through a batch engine. Everything that could otherwise confound a
comparison between mechanisms --- the asset, the participant
population, the information environment, the time period, the
volatility regime --- is held fixed by construction, because it is
literally the same event stream. The matching engine is the only
treatment variable.

\section{Simulation Environment}
\label{sec:environment}

\subsection{Shared Type System and Matching Interface}

Both engines are implemented in Rust and share one order, trade, and
book type system, and both implement one common matching interface.
This is not a stylistic choice. If the two engines had separate order
representations or separate event loops, any measured difference could
originate in the representation rather than in the mechanism. A shared
interface makes the engine substitutable at a single point in the
harness and makes the claim ``the engine is the only difference''
verifiable rather than asserted.

\begin{lstlisting}[caption={The shared matching interface. Both engines
implement the same trait, so the harness is generic over the
mechanism.},label={lst:trait}]
pub trait MatchingEngine {
    /// Submit, modify or cancel. Returns events generated
    /// immediately (empty for the FBA during accumulation).
    fn on_event(&mut self, ev: &InboundEvent) -> Vec<EngineEvent>;

    /// Advance the engine clock. The FBA clears here when a
    /// batch boundary is crossed; the CDA is a no-op.
    fn on_clock(&mut self, now: Nanos) -> Vec<EngineEvent>;

    /// Snapshot of the current book state for metric computation.
    fn book_state(&self, pair: PairId) -> BookState;

    /// Engine identity and configuration, written to the log
    /// header for reproducibility.
    fn config(&self) -> EngineConfig;
}
\end{lstlisting}

\subsection{The Continuous Engine (CDA)}

The continuous engine maintains a price-time priority limit order book
per asset pair. Each incoming order is matched immediately against the
opposite side of the book, executing at the resting maker price, and any
unfilled remainder rests according to its time-in-force. Cancellations
and modifications are applied in arrival order; a modification that
increases size or improves price loses time priority, following the
convention of the venue whose data is replayed.

This engine is the control condition. It is not intended to be a novel
contribution but a faithful reference implementation, and
\Cref{sec:validation} describes how its fidelity is checked against the
trades that the source venue actually produced from the same flow.

\subsection{The Batch Engine (FBA)}

The batch engine buffers all events arriving within an interval $\tau$
and clears at the end of the interval. Clearing proceeds in the stages
of \Cref{fig:pipeline}: the accumulated buy and sell schedules are
constructed as in \eqref{eq:demandsupply}, the volume-maximising price
is found by a search over the price grid as in \eqref{eq:clearing}, ties
are resolved by the documented cascade, inframarginal orders are filled
completely, and the marginal side is rationed by the configured rule
from \Cref{sec:rationing}.

Volume that cannot be matched at $\clr$ is not executed. It is rolled
over into the next batch unless its time-in-force says otherwise. No
external liquidity source, such as an automated market maker, is used to
fill the residual. This is what keeps the comparison strictly
like-for-like: both engines can only execute against the order flow that
was actually submitted, so neither has access to liquidity the other
lacks.

The engine is configurable along the dimensions derived in
\Cref{ch:design}: $\tau$, fixed or randomised boundary, rationing rule,
tie-break cascade, and residual policy. \Cref{sec:experiments} states
which of these are varied.

\subsection{The Replay Harness}

A shared harness reads the event stream, dispatches each event to the
engine under test in timestamp order, triggers clock ticks so that the
FBA's batch boundaries fall where the configuration says they should,
and writes every order, trade, and book state transition to a structured
log for offline metric computation. Metrics are never computed inside
the engine; they are computed from the log by a separate analysis
pipeline, so that both engines are measured by exactly the same code.

Two properties of the harness matter for the validity of the results.
It is deterministic: given the same input stream and the same
configuration, including the seed for any randomised rule, it produces
byte-identical logs. And it separates simulation time from wall-clock
time, so that engine throughput and clearing latency can be measured
without those measurements affecting the simulated market.

\section{Data Source and Data Levels}
\label{sec:data}

\subsection{The Hyperliquid Level~4 Dataset}

The simulation is fed with the publicly released Hyperliquid Level~4
order book dataset of \citet{albers2026openbook}, captured directly at a
non-validating Hyperliquid node and published on Zenodo. Beyond a
standard L4 feed it also records rejected orders, failed cancellations,
both counterparties to every trade, and post-trade inventory, all at
nanosecond precision. Because it is captured at the node rather than
sampled from periodic snapshots, no book reconstruction step is
required, which removes a substantial source of error that affects
studies working from vendor feeds.

Three properties make this dataset the right choice here. Order-by-order
granularity is the necessary condition for replaying historical flow
through a different matching engine at all. Participant attribution is
what makes the allocation and per-participant metrics of RQ2.3
computable. And nanosecond timestamps are what make interval boundary
placement meaningful at the millisecond scale of the batch intervals
tested.

\subsection{The Data-Level Framework}

\Cref{tab:levels} sets out the standard hierarchy of order book data and
what each level enables. The framework serves two purposes in this
thesis. It scopes the metric catalogue, since every metric in
\Cref{sec:metrics} is tagged with the minimum level it requires. And it
makes the contribution of the dataset explicit by showing which metrics
would simply be unavailable at a lower level.

\begin{table}[htbp]
\centering
\caption[Order book data levels]{The order book data-level hierarchy
and what each level enables. This thesis works with L4 throughout.}
\label{tab:levels}
\small
\begin{tabularx}{\textwidth}{@{}p{1.0cm} p{4.2cm} L@{}}
\toprule
\textbf{Level} & \textbf{Content} & \textbf{What it enables}\\
\midrule
L1 & Best bid and offer, last trade & Quoted spread, midpoint series,
realized volatility, effective spread\\
L2 & Aggregated depth per price level (the ladder) & Depth at best and
within $x$ bps, depth curve shape, book imbalance, slippage of a
hypothetical order\\
L3 & Order-by-order feed: every add, modify, cancel and execution with
an order ID and timestamp & Exact book reconstruction, queue position,
time priority effects, cancellation and order-to-trade ratios,
arrival-time distribution inside an interval\\
L4 & L3 plus participant attribution (wallet or account identifier) per
order and per fill & Per-participant metrics: maker inventory and
markout PnL, adverse selection by counterparty type, fill distribution
across participants under a rationing rule, order size inflation, flow
classification into benign and toxic\\
\bottomrule
\end{tabularx}
\end{table}

\subsection{What the Simulation Observes and What It Proxies}

Because the engines are the author's own, several quantities that are
unobservable in field data become directly observable here: every
submitted limit price including those of orders that never execute, the
exact rationing outcome order by order, the internal clearing state
including the full demand and supply schedules at every candidate price,
and the volume left unexecuted after each batch. Field studies must
infer these or do without them.

What remains unobservable is proxied or modelled instead, and stating
this explicitly is part of the design:

\begin{itemize}
  \item \textbf{The true fundamental value} is proxied by Hyperliquid's
  own oracle and mark price, which is used as the external reference in
  the pricing error metric.
  \item \textbf{Hidden liquidity} that was never submitted as a visible
  order cannot be recovered and is absent from both engines equally.
  \item \textbf{True participant latency} is not observable; the
  timestamps record when the node saw the message, not when the
  participant decided to send it.
\end{itemize}

The third of these bears on interpretation. Because arrival times are
those produced under Hyperliquid's own continuous engine, the position
of an order relative to a simulated batch boundary is an artefact of
where the boundary was placed, not a choice the participant made. The
consequences are drawn out in \Cref{sec:testable}.

\section{Experimental Design and Comparability Protocol}
\label{sec:experiments}

\subsection{Treatment and Controls}

The treatment is the matching engine. The comparability protocol
enforces that nothing else differs:

\begin{enumerate}
  \item \textbf{Identical input.} Both runs consume the same event
  stream, byte for byte, in the same order.
  \item \textbf{Identical initial state.} Both engines start from an
  empty book and are warmed up over the same lead-in period, which is
  excluded from all metric computation.
  \item \textbf{Identical metric code.} Metrics are computed offline
  from the logs by one implementation shared by both engines.
  \item \textbf{Identical treatment of edge cases.} Self-matching,
  rejected orders, and failed cancellations are handled by the same
  policy in both engines, documented in the configuration header of
  every log.
  \item \textbf{Determinism.} Every run is reproducible from its
  configuration and seed.
\end{enumerate}

\subsection{Sample Selection}

The sample is drawn along three dimensions so that results are not an
artefact of one instrument or one market state:

\begin{itemize}
  \item \textbf{Instruments:} a liquid major pair and a less liquid pair,
  since \Cref{sec:interval} predicts that interval length interacts with
  the volume accumulating per batch.
  \item \textbf{Sessions:} several full trading days, plus at least one
  high-volatility episode, since the value of removing latency
  arbitrage should be state-dependent.
  \item \textbf{Time of day:} intervals of high and low activity, to
  separate mechanism effects from activity effects.
\end{itemize}

\subsection{Parameter Sweeps}

Two sweeps operationalise design blocks from \Cref{ch:design}:

\begin{itemize}
  \item \textbf{Interval sweep.} $\tau \in \{1, 10, 50, 100, 250, 500,
  1000\}$ milliseconds, spanning the range from below to well above the
  latency dispersion the mechanism is meant to neutralise. This sweep
  tests prediction~P8 and traces the delay-versus-thickness trade-off of
  \Cref{sec:interval}.
  \item \textbf{Rationing sweep.} Price-time within batch, pro rata,
  size priority, and random allocation, holding $\tau$ fixed at the
  baseline. Because the flow is historical, this sweep measures the
  mechanical redistribution each rule produces on a fixed order set, not
  a behavioural response.
\end{itemize}

The baseline configuration against which the CDA is compared is
$\tau = 100$\,ms, fixed boundary, pro rata rationing, roll-over
residual, tie-break cascade as in \Cref{sec:clearing}.

\subsection{Statistical Treatment}

The comparison is paired by construction: each interval of the source
data produces one CDA observation and one FBA observation from the same
flow. Differences are therefore evaluated as paired differences per
interval rather than as two independent samples, which removes the
common variation driven by the market state and leaves the mechanism
effect. Because microstructure series are strongly autocorrelated,
inference uses standard errors robust to autocorrelation, and results
are reported per instrument and per session as well as pooled, so that
a single unusual session cannot drive a pooled conclusion.

\section{Metric Catalogue}
\label{sec:metrics}

Notation: $a_t$ and $b_t$ are the best ask and bid, $\midp = (a_t+b_t)/2$
the midpoint, $p_k$ the execution price of trade $k$, $q_k$ its size,
$D_k \in \{+1,-1\}$ its sign (buyer- or seller-initiated), $\pi_k$ the
submitted limit price, $\clr$ the uniform clearing price of a batch, and
$\Delta$ a markout horizon. The \emph{Data} column gives the minimum
data level required.

\begin{table}[htbp]
\centering
\caption[Metrics for RQ2.1]{RQ2.1 --- liquidity and transaction cost
metrics.}
\label{tab:metrics1}
\small
\begin{tabularx}{\textwidth}{@{}p{3.1cm} L p{1.6cm}@{}}
\toprule
\textbf{Metric} & \textbf{Definition and measurement} & \textbf{Data}\\
\midrule
Quoted spread & $(a_t-b_t)/\midp$, time-weighted per interval; for the
FBA the implied spread between the marginal unexecuted buy and sell of
the batch & L1/L2\\
Depth at best & Displayed volume at the best bid and ask, time-weighted;
for the FBA the volume at $\clr$ in the batch schedules & L2\\
Depth within $x$ bps & Cumulative volume within $x$ basis points of
$\midp$ (or of $\clr$), for several $x$ & L2\\
Book imbalance & $(V^{bid}-V^{ask})/(V^{bid}+V^{ask})$ at the top $n$
levels & L2\\
Effective spread & $2D_k(p_k-m_k)/m_k$, volume-weighted; $m_k$ is the
pre-trade midpoint, for the FBA the midpoint at the start of the batch
& L1 + trades\\
Realized spread & $2D_k(p_k-m_{k+\Delta})/m_k$ for
$\Delta \in \{1,5,30\}$ seconds; the part of the spread the liquidity
provider keeps & L1 + trades\\
Price impact & Effective spread minus realized spread; the adverse
selection component & L1 + trades\\
Amihud illiquidity & $|r_t|/\text{volume}_t$ averaged over intervals &
L1 + trades\\
\bottomrule
\end{tabularx}
\end{table}

\begin{table}[htbp]
\centering
\caption[Metrics for RQ2.2]{RQ2.2 --- price discovery and market quality
metrics.}
\label{tab:metrics2}
\small
\begin{tabularx}{\textwidth}{@{}p{3.1cm} L p{1.6cm}@{}}
\toprule
\textbf{Metric} & \textbf{Definition and measurement} & \textbf{Data}\\
\midrule
Realized volatility & Standard deviation of interval midpoint or
clearing price returns & L1\\
Intra-interval price dispersion & Standard deviation of execution prices
within one interval; zero by construction for the FBA, strictly positive
for the CDA, and therefore a direct measure of the uniform-price
property & L1 + trades\\
Pricing error & $|\clr-p^{ref}|$ or $|\midp-p^{ref}|$ against an
external reference price (Hyperliquid's oracle/mark price), in bps &
L1 + ref\\
\bottomrule
\end{tabularx}
\end{table}

\begin{table}[htbp]
\centering
\caption[Metrics for RQ2.3]{RQ2.3 --- execution, allocation and
order-arrival timing metrics, plus engine controls.}
\label{tab:metrics3}
\small
\begin{tabularx}{\textwidth}{@{}p{3.1cm} L p{1.6cm}@{}}
\toprule
\textbf{Metric} & \textbf{Definition and measurement} & \textbf{Data}\\
\midrule
Executed volume & Matched quantity and notional per interval; number of
trades & L1 + trades\\
Fill rate & Filled quantity divided by submitted quantity, overall and
by participant class & L3/L4\\
Time to execution & Elapsed time from submission to fill; for the FBA
bounded below by the time to the next batch boundary, which quantifies
the delay cost of batching & L3\\
Trader surplus & $\sum_k |\pi_k - \clr|\, q_k$ over executed orders: the
gain relative to the submitted limit price & L3/L4\\
Order size inflation & Ratio of submitted to filled size per participant
over time; the empirical signature of order stuffing under pro rata &
L4\\
Order-to-trade ratio & Submitted messages per executed trade;
cancellation rate per interval; measures message traffic and quote
flickering & L3\\
Boundary concentration & Share of order arrivals in the final $x\%$ of
the batch interval & L3\\
Throughput & Orders processed per second under identical load &
control\\
Clearing latency & Wall-clock time per batch clearing computation and
per continuous match & control\\
Unexecuted residual & Share of batch volume on the heavier side that
cannot be matched at $\clr$ and remains unexecuted (rolled over or
cancelled, per time-in-force) & control\\
\bottomrule
\end{tabularx}
\end{table}

\subsection{Interpretive Status of the RQ2.3 Metrics}
\label{sec:testable}

Three metrics in \Cref{tab:metrics3} require an interpretive caveat that
is easy to state and important not to forget when reading
\Cref{ch:results}.

Order size inflation, boundary concentration, and to a lesser extent
trader surplus are all quantities whose theoretical interest lies in
their being \emph{responses} to the mechanism. Order size inflation is
interesting because participants inflate size when facing pro rata
allocation. Boundary concentration is interesting because participants
cluster submissions near a predictable gate.

The orders replayed here were submitted under Hyperliquid's own
continuous engine. Their sizes were not chosen in response to the
rationing rule under test, and their arrival times were not chosen in
response to a batch boundary that did not exist. What the simulation
measures is therefore the mechanical effect of the rule on fixed
historical flow: how the pro rata rule would have allocated these
particular orders, and how this particular arrival-time distribution
happens to fall relative to a boundary chosen for the simulation. These
are meaningful quantities --- they establish the baseline against which
a behavioural response would have to be measured --- but they are not
evidence about how participants would behave. Predictions P5, P6, and P7
of \Cref{tab:predictions} are therefore out of reach of this design, and
\Cref{sec:futurework} discusses the agent-based extension that would
reach them.

The throughput and clearing latency metrics serve a different purpose.
They are not results about market quality but a control: if the two
engines differed substantially in processing performance, a measured
difference in market outcomes could reflect implementation rather than
mechanism. Reporting them is how that alternative explanation is ruled
out.

\section{Validation of the Harness}
\label{sec:validation}

A simulation result is only as credible as the machinery that produced
it. Four layers of validation are applied.

\paragraph{Invariant checks.} Every engine run asserts a set of
invariants continuously: quantity conservation (executed buy volume
equals executed sell volume in every batch and every match), price
bounds (no buyer pays more than its limit, no seller receives less), no
negative or duplicated fills, and monotone timestamps. A violated
invariant aborts the run rather than producing a quietly wrong log.

\paragraph{Unit tests against hand-computed cases.} The clearing routine
is tested against small books whose clearing price and allocation can be
computed by hand, including the awkward cases: an empty book, no
crossing orders, a unique clearing price, a clearing interval requiring
each stage of the tie-break cascade, an exactly balanced batch requiring
no rationing, and a batch in which a single order is pivotal.

\paragraph{Replay fidelity of the CDA.} The continuous engine is
replayed against the same flow that Hyperliquid's own engine processed,
and its output is compared with the trades the venue actually produced.
Perfect agreement is not expected, since matching details, self-trade
prevention, and message handling at the node are not fully specified in
the public documentation, but the agreement rate is reported as a
measure of how faithful the control condition is. A low agreement rate
would undermine the comparison, since the CDA would then not be the
mechanism that generated the flow.

\paragraph{Degenerate-parameter checks.} As $\tau \to 0$ the FBA should
converge toward the CDA in cleared volume and price path, since a batch
containing at most one event is a continuous match. Verifying this
convergence numerically is a strong end-to-end test of the batch
engine: it exercises the whole pipeline and has a known answer.
